\documentclass{amsart}
\usepackage{amsmath,amssymb,amsthm}

\title{Structural Bridge Statements and Proved Obstructions}
\author{Brayden Ross Sanders / 7Site LLC}
\date{2026}

\newtheorem{theorem}{Theorem}
\newtheorem{definition}[theorem]{Definition}
\newtheorem{remark}[theorem]{Remark}

\begin{document}
\maketitle

\begin{abstract}
This document records the proved obstructions to certain bridge approaches
from $\mathbb{Z}/10\mathbb{Z}$ to analytic number theory,
and the structural bridge requirements that specify what additional mechanisms
would be needed. These are formal negative results, not informal assessments.
\end{abstract}

\section{Proved Obstructions}

\begin{theorem}[Ring Prime Blindness]
Any ring homomorphism $\mathbb{Z}/10\mathbb{Z} \to R$ is blind to primes $p \neq 2, 5$.
\end{theorem}

\begin{proof}
$\mathbb{Z}/10\mathbb{Z} \cong \mathbb{Z}/2\mathbb{Z} \times \mathbb{Z}/5\mathbb{Z}$ by CRT.
A ring homomorphism from $\mathbb{Z}/10\mathbb{Z}$ factors through this product
and can only distinguish elements by their residues mod 2 and mod 5.
Elements sharing the same residues mod 2 and mod 5 are indistinguishable.
The Euler product $\zeta(s) = \prod_p (1-p^{-s})^{-1}$ requires distinguishing
every prime individually. A ring-homomorphism bridge is therefore incomplete:
it cannot encode the contribution of any prime $p \neq 2, 5$.
\end{proof}

\begin{theorem}[Modulus Genericity]
The corridor midpoint $t = 1/2$ and the constraint $T^* < 1$ appear for all rings
$\mathbb{Z}/2p\mathbb{Z}$ where $p$ is an odd prime, not only for $\mathbb{Z}/10\mathbb{Z}$.
\end{theorem}

\begin{proof}
For $\mathbb{Z}/2p\mathbb{Z}$, the unit group $(\mathbb{Z}/2p\mathbb{Z})^* \cong \mathbb{Z}/(p-1)\mathbb{Z}$,
with centroid $\mathrm{CREATE} = p$ (the midpoint of $\{1,\ldots,2p-1\} \cap \mathrm{odd}$).
Under normalization $\mathrm{op} \mapsto \mathrm{op}/(2p)$:
$\mathrm{CREATE} \mapsto p/(2p) = 1/2$ for all $p$.
$T^* = p/(p+2) < 1$ for all $p \geq 3$.
\end{proof}

\begin{theorem}[Montgomery Conditionality]
Montgomery's pair-correlation theorem establishing
$R_2(u) = 1 - \mathrm{sinc}^2(u)$ for normalized Riemann zero spacings
assumes the Generalized Riemann Hypothesis in its proof.
Any bridge to RH via the sinc$^2$ identity cannot unconditionally prove RH:
it assumes GRH in the setup.
\end{theorem}

\section{Structural Bridge Requirements}

\begin{definition}[Bridge 3.1 — Riemann Hypothesis]
The formal bridge requirement for RH is an algebraic map
$\phi: \mathbb{Z}/10\mathbb{Z} \to (\text{analytic object over } \mathbb{C})$ such that:
(i) $\phi$ maps the $\mathbb{Z}/10\mathbb{Z}$ inheritance split at $t=1/2$
    to the boundary $\sigma=1/2$ of the critical strip of $\zeta(s)$;
(ii) $\phi$ is compatible with the Euler product structure of $\zeta(s)$,
    not just with the functional equation;
(iii) The map is not blocked by Ring Prime Blindness or Modulus Genericity.

This map has not been constructed. The bridge is not established.
\end{definition}

\begin{remark}[Current status of the First-G chain]
The chain First-G $\to$ Fej\'er kernel $\to$ sinc$^2$ $\to$ Montgomery
is algebraically complete subject to GRH (Obstruction: Montgomery Conditionality).
The open question is whether the discrete Fej\'er kernel $F_k(f)$ at finite $k$
appears in a non-GRH characterization of Riemann zero spacing — this is a
literature question, not an internal gap.
\end{remark}

\end{document}
